\documentclass[11pt,a4paper]{article}
\usepackage[margin=2.4cm]{geometry}
\usepackage{amsmath,amssymb,amsthm,mathtools}
\usepackage[protrusion=true,expansion=false]{microtype}
\usepackage{booktabs,enumitem}
\usepackage{tikz}
\usetikzlibrary{arrows.meta,positioning,calc,fit,backgrounds}
\definecolor{tierspend}{RGB}{200,60,60}
\definecolor{tierfull}{RGB}{210,140,40}
\definecolor{tierin}{RGB}{60,130,180}
\definecolor{tierrecv}{RGB}{70,150,90}
\usepackage[colorlinks=true,linkcolor=blue!50!black,citecolor=blue!50!black,urlcolor=blue!50!black]{hyperref}

\emergencystretch=3em
\hbadness=10000
\hfuzz=2pt

\newtheorem{theorem}{Theorem}[section]
\newtheorem{lemma}[theorem]{Lemma}
\newtheorem{proposition}[theorem]{Proposition}
\newtheorem{corollary}[theorem]{Corollary}
\theoremstyle{definition}
\newtheorem{definition}[theorem]{Definition}
\newtheorem{construction}[theorem]{Construction}
\newtheorem{assumption}[theorem]{Assumption}
\newtheorem{remark}[theorem]{Remark}
\newtheorem{example}[theorem]{Example}

\newcommand{\F}{\mathbb{F}}
\newcommand{\Z}{\mathbb{Z}}
\newcommand{\N}{\mathbb{N}}
\newcommand{\G}{\mathbb{G}}
\newcommand{\E}{\mathbb{E}}
\newcommand{\PP}{\mathbb{P}}
\renewcommand{\Pr}{\mathrm{Pr}}
\newcommand{\Fp}{\F_p}
\newcommand{\Fq}{\F_q}
\newcommand{\ip}[2]{\langle #1,\, #2 \rangle}
\newcommand{\Comm}{\mathsf{Commit}}
\newcommand{\Open}{\mathsf{Open}}
\newcommand{\Setup}{\mathsf{Setup}}
\newcommand{\Verify}{\mathsf{Verify}}
\newcommand{\negl}{\mathsf{negl}}
\newcommand{\poly}{\mathsf{poly}}
\newcommand{\PPT}{\mathsf{PPT}}
\newcommand{\Adv}{\mathcal{A}}
\newcommand{\Prov}{\mathcal{P}}
\newcommand{\Ver}{\mathcal{V}}
\newcommand{\Sim}{\mathcal{S}}
\newcommand{\Ext}{\mathcal{E}}
\newcommand{\Rel}{\mathcal{R}}
\newcommand{\Lang}{\mathcal{L}}
\newcommand{\nfsym}{\mathsf{nf}}
\newcommand{\cmsym}{\mathsf{cm}}
\newcommand{\cmx}{\mathsf{cmx}}
\newcommand{\cvsym}{\mathsf{cv}}
\newcommand{\rk}{\mathsf{rk}}
\newcommand{\Pallas}{\ensuremath{\mathsf{Pallas}}}
\newcommand{\Vesta}{\ensuremath{\mathsf{Vesta}}}
\newcommand{\Ep}{\ensuremath{\mathcal{E}}}
\newcommand{\NoteCommit}{\ensuremath{\mathsf{NoteCommit}}}
\newcommand{\Sinsemilla}{\ensuremath{\mathsf{Sinsemilla}}}
\newcommand{\ShortCommit}{\ensuremath{\mathsf{ShortCommit}}}
\newcommand{\Extract}{\ensuremath{\mathsf{Extract}}}
\newcommand{\Acc}{\ensuremath{\mathsf{Acc}}}
\newcommand{\GroupHash}{\ensuremath{\mathsf{GroupHash}}}
\newcommand{\ask}{\ensuremath{\mathsf{ask}}}
\newcommand{\aksym}{\ensuremath{\mathsf{ak}}}
\newcommand{\nksym}{\ensuremath{\mathsf{nk}}}
\newcommand{\ivk}{\ensuremath{\mathsf{ivk}}}
\newcommand{\fvk}{\ensuremath{\mathsf{fvk}}}
\newcommand{\ovk}{\ensuremath{\mathsf{ovk}}}
\newcommand{\dk}{\ensuremath{\mathsf{dk}}}
\newcommand{\pkd}{\ensuremath{\mathsf{pk_d}}}
\newcommand{\gd}{\ensuremath{\mathsf{g_d}}}
\newcommand{\rcv}{\ensuremath{\mathsf{rcv}}}
\newcommand{\rcm}{\ensuremath{\mathsf{rcm}}}
\newcommand{\rivk}{\ensuremath{\mathsf{r_{ivk}}}}
\newcommand{\rsk}{\ensuremath{\mathsf{rsk}}}

\title{\textbf{\Huge Halo 2 Guide}\\[6pt]\large The Halo 2 proof system, from arithmetisation to recursion}
\author{m@rek.onl}
\date{}
\begin{document}
\maketitle
\thispagestyle{empty}
\begin{abstract}
This is the third volume of a four-part series --- \emph{Math Guide}, \emph{Crypto Guide}, \emph{Halo~2 Guide}, \emph{Orchard Guide}. Assuming the mathematics of the \emph{Math Guide} and the cryptographic primitives of the \emph{Crypto Guide}, it develops the Halo~2 proof system: the polynomial-IOP design pattern, a rigorous account of knowledge soundness in the algebraic group model, proof-carrying data and accumulation schemes, and finally Halo~2 itself --- PLONKish arithmetisation, the permutation and lookup arguments, the inner-product polynomial commitment and the Halo accumulation trick, and the complete protocol.
\end{abstract}
\clearpage
\tableofcontents
\clearpage

\section{Introduction}\label{sec:intro}
This is the \emph{Halo~2 Guide}, the third volume of a four-part development of the cryptography
behind Zcash's Orchard protocol. The first two volumes --- the \emph{Math Guide} and the
\emph{Crypto Guide} --- constructed, from the ground up, the mathematics (finite fields, elliptic curves,
polynomials and evaluation domains, probability) and the cryptographic primitives (hash functions,
commitments, $\Sigma$-protocols and the Fiat--Shamir transform, zero knowledge, polynomial
commitment schemes) that this volume assumes. The present volume assembles those ingredients into a
\emph{proof system}.

The development proceeds in four steps. First, we make the polynomial-IOP design pattern precise: the
recipe ``encode a computation as polynomial identities over an evaluation domain, then compile to a
succinct argument with a polynomial commitment.'' Next, we treat knowledge soundness rigorously,
including the algebraic group model and the tree-of-transcripts extractor that the inner-product
argument requires. We then develop proof-carrying data, incrementally-verifiable computation, and
the accumulation schemes that make recursion practical. Finally, we assemble Halo~2 itself:
PLONKish arithmetisation, the permutation and lookup arguments, the inner-product polynomial
commitment with a worked example and its accumulation (the ``Halo trick''), and the complete
seven-phase protocol. The \emph{Orchard Guide} subsequently uses this proof system to construct a shielded
payment protocol.

\section{Polynomial IOPs and the SNARK design pattern}\label{sec:piop}

The modern SNARK design pattern factors the protocol into two clean
layers: a \emph{polynomial interactive oracle proof} (PIOP) for the
relation, and a polynomial commitment scheme (the \emph{Crypto Guide}) that
realises the PIOP's oracles. This factoring, due in its current form to
Bünz, Fisch, Szepieniec, and others, has substantially clarified the
nature of SNARKs.

\subsection{The polynomial IOP, formally}

\begin{definition}[Polynomial IOP]\label{def:piop}
A \emph{polynomial interactive oracle proof} for a relation $\Rel$ is a
public-coin interactive protocol where in each round:
\begin{itemize}[leftmargin=*]
\item The prover sends one or more \emph{polynomial oracles} $p_i \in \F[X]$
(each of bounded degree).
\item The verifier sends a random challenge $c \in \F$.
\end{itemize}
At the end of the protocol, the verifier may query each oracle at a small
number of points, learning the polynomial's evaluations at those points.
The verifier's final accept/reject is a deterministic function of the
challenges, evaluations, and the public input $x$.

We define completeness, soundness, and knowledge soundness as usual.
\end{definition}

A PIOP differs from a vanilla interactive argument in one essential respect:
the verifier does not read the prover's messages in full --- it queries
only a few evaluations of each polynomial oracle. This is what renders the
verifier ``succinct'' even when the polynomials themselves have large degree.

\subsection{Compiling a PIOP into a SNARK}

\begin{construction}[PIOP-to-SNARK compiler]\label{const:piop-snark}
Given a PIOP $\Pi$ for $\Rel$, a PCS scheme $(\Setup, \Comm, \Open, \Verify)$
for $\F[X]$, and a random oracle $H$, the compiler produces the SNARK:
\begin{enumerate}[leftmargin=*]
\item Setup: run the PCS setup to obtain $\mathsf{pp}$.
\item Prover: simulate $\Pi$, replacing each polynomial oracle $p_i$ by its
PCS commitment $C_i$, and each verifier-challenge / oracle-query round by
the corresponding PCS opening. Use Fiat--Shamir (the \emph{Crypto Guide}) to derive
challenges from the transcript.
\item Verifier: parse the transcript, recompute Fiat--Shamir challenges,
verify each PCS opening, and run the PIOP's final accept/reject check.
\end{enumerate}
\end{construction}

\begin{theorem}[Soundness preservation]\label{thm:piop-snark}
If $\Pi$ has knowledge soundness $\kappa_{\Pi}$ and the PCS has knowledge
binding $\kappa_{\mathrm{PCS}}$, the compiled SNARK has knowledge soundness
$O(\kappa_\Pi + Q \cdot \kappa_{\mathrm{PCS}})$ in the random oracle model,
where $Q$ is the adversary's RO-query budget.
\end{theorem}

This factoring is of considerable utility: research on PIOPs (PLONK, Marlin,
Spartan, HyperPlonk, $\dots$) proceeds independently of research on
polynomial commitment schemes, and one may combine any compatible pair to
produce a new SNARK with mixed-and-matched properties. This document
develops the one PCS that Halo 2 uses --- the inner-product argument
(\S\ref{sec:ipa-detail}) --- and treats the PIOP/PCS split only insofar as
it clarifies the Halo 2 construction.

\subsection{A worked example: a PIOP for univariate evaluation}\label{ex:univariate-piop}

To render the abstraction concrete, consider a minimal PIOP for the simplest
non-trivial relation: ``the prover knows a polynomial $p$ of degree $\le d$
such that $p(z) = y$ for a public $(z, y)$.''

\begin{construction}[PIOP for univariate evaluation]\label{const:univ-piop}
On input $(z, y)$ and prover witness $p$:
\begin{enumerate}[leftmargin=*]
\item Prover sends the polynomial oracle $p$ and the quotient oracle
$q := (p - y)/(X - z)$.
\item Verifier samples a random challenge $r \in \F$.
\item Verifier queries $p$ and $q$ at $r$.
\item Verifier accepts iff $p(r) - y = q(r) \cdot (r - z)$.
\end{enumerate}
\end{construction}

Completeness is immediate ($q$ is a polynomial since $p(z) - y = 0$).
Soundness follows from Schwartz--Zippel: if $(X - z) \nmid (p(X) - y)$,
then $p(X) - y - q(X)(X - z)$ is a non-zero polynomial of degree $\le d$,
and its random evaluation is zero with probability $\le d/|\F|$.

Compiling this PIOP with the IPA-PCS, applying Fiat--Shamir, yields a
non-interactive SNARK for the same relation. The PCS handles polynomial
hiding and binding; the PIOP handles the relation logic. This separation
is precisely the modular SNARK design pattern.

\subsection{PLONK as a PIOP}

PLONK (\S\ref{sec:plonkish}) is a more elaborate PIOP. Its
oracles are:
\begin{itemize}[leftmargin=*]
\item One advice polynomial per advice column.
\item A permutation running-product polynomial $Z_{\mathrm{perm}}$.
\item One lookup running-product polynomial per lookup.
\item A quotient polynomial $h$ (split into chunks).
\end{itemize}
Its challenges are $\beta, \gamma$ (for the permutation argument),
$\theta$ (for lookup column compression), $y$ (for the gate-equation
combination), $x$ (the final evaluation point). Its final check is the
PLONK gate equation $C(x) = h(x) \cdot Z_H(x)$, which the verifier checks from
the evaluations of the oracles at $x$ (and possibly $\omega x$).

This is precisely the PIOP that phases 1--6 of the Halo 2
protocol of \S\ref{sec:halo2-full} compile; phase 7 is the IPA-PCS
opening of the combined polynomial.

\section{A rigorous treatment of knowledge soundness}\label{sec:ks}

We now develop knowledge soundness with the care needed for arguments
about modern SNARKs. The treatment combines the rewinding-extractor
framework (the \emph{Crypto Guide}) with the algebraic group model.

\subsection{Definitional refinements}

The bare definition of the \emph{Crypto Guide} suffices for $\Sigma$-protocols
but is awkward for compound systems. We adopt the following standard refinement.

\begin{definition}[Witness-extended emulation]\label{def:wee}
An interactive argument $(\Prov, \Ver)$ has \emph{witness-extended
emulation} if there is a $\PPT$ emulator $\Ext$ such that for every $\PPT$
prover $\Prov^*$ and every input $x$:
\begin{itemize}[leftmargin=*]
\item Run $\Prov^*$ honestly with $\Ver$; let $\tau$ be the transcript and
$b$ the verifier's accept/reject bit.
\item Run $\Ext^{\Prov^*}(x)$; it outputs $(\tau', b', w)$.
\end{itemize}
The distributions of $(\tau, b)$ and $(\tau', b')$ are computationally
indistinguishable, and moreover whenever $b' = 1$, $(x, w) \in \Rel$ except
with negligible probability.
\end{definition}

WEE captures simultaneously two properties: ``the extractor produces a witness whenever
the protocol accepts'' and ``the extractor's external behaviour resembles
a normal protocol execution.'' It is the appropriate notion for arguments
embedded in larger protocols (e.g., for batched proofs or for verifying
Halo 2 proofs within a Halo 2 verifier circuit).

\subsection{The algebraic group model}

The IPA's knowledge-soundness proof proceeds through the tree extractor of
\S\ref{sec:ipa-detail}, with knowledge error $O(d / |\F|)$. The number of
prover invocations the extractor makes is $O(d)$, but converting those
invocations into a standard-model reduction to DL via iterated forking is
loose: each of the $\log d$ rounds contributes a multiplicative blow-up,
so the standard-model extractor runs in quasi-polynomial time
$d^{O(\log d)}$ rather than strict polynomial time. The algebraic group
model removes this loss.

The \emph{algebraic group model} (AGM; Fuchsbauer, Kiltz, Loss; CRYPTO
2018) provides a tighter analysis by restricting the adversary's group
operations.

\begin{definition}[Algebraic adversary]\label{def:algebraic}
An \emph{algebraic adversary} against a group $\G$ is a $\PPT$ algorithm
that, whenever it outputs a group element $P \in \G$, additionally outputs
an explicit linear combination
$P = \sum_i [a_i] Q_i$ expressing $P$ as a combination of previously-seen
group elements $Q_i$.
\end{definition}

The AGM is intermediate between the standard model (no restrictions on the
adversary) and the generic group model (the adversary can access group
operations only via an oracle). Many real adversaries are algebraic --- they
compute group elements via known scalar multiplications and additions ---
so AGM analyses are often sharper than standard-model analyses while
remaining heuristically conservative.

For Halo 2, BCMS20 conduct the accumulation analysis in the AGM. The resulting
extractor is polynomial in the round count $k = \log d$, neither exponential
nor quasi-polynomial: each algebraic adversary's group-element outputs
come with explicit linear combinations, so the extractor reads
off the witness from the combinations.

\subsection{Tree-extractor analysis}

To render the tree extractor concrete, consider the key step at each round of the
IPA halving (see \S\ref{sec:ipa-detail}):

\begin{lemma}[Single-round extraction]\label{lem:single-round}
In a single round of the IPA at depth $j$, given three accepting
sub-transcripts that share the round's first message $(L_j, R_j)$ but
have three distinct challenges $u_j$, the extractor recovers the round-$j$
witness halves $\vec{a}_{\mathrm{lo}}^{(j)}, \vec{a}_{\mathrm{hi}}^{(j)}$
together with the consistency of the $L_j, R_j$ cross-terms via $3 \times 3$
Vandermonde inversion (rows $(u_j^{-2}, 1, u_j^2)$), since the folded
relation is quadratic in $u_j$.
\end{lemma}

\begin{proof}[Sketch]
The three accepting transcripts give three accepting verifier equations.
Because the folded relation is quadratic in $u_j$ (the three monomials
$u_j^{-2}, 1, u_j^2$), treating these as three linear equations in three
unknowns (the cross-terms and the folded commitment), the system has a
unique solution because the $3 \times 3$ Vandermonde determinant with rows
$(u_j^{-2}, 1, u_j^2)$ is non-zero when the three challenges are distinct
(and $u_j \ne \pm u_j'$, etc.). Excluding the bad challenge collisions adds
a factor $O(1/|\F|)$ to the knowledge error.
\end{proof}

The full extraction requires three transcripts at each round, giving a tree of
$3^k$ leaves and $O(d)$ prover invocations. In the AGM, each prover
invocation is direct (the extractor reads the algebraic representation);
in the standard model with forking, each invocation costs $1/p^*$ expected
queries to the prover, yielding total cost $O(d \cdot \mathsf{poly}(\lambda) / p^*)$.

\subsection{Concrete versus asymptotic security}

The literature on knowledge soundness contains numerous tightness gaps. The
more important ones include:
\begin{itemize}[leftmargin=*]
\item Bulletproofs' standard-model knowledge-soundness reduction loses a
factor exponential in the round count $\log d$ (equivalently,
quasi-polynomial $d^{O(\log d)}$ in the degree bound $d$).
\item Halo 2's deployed parameters assume the AGM and/or the random oracle
model heuristically.
\item PLONK's original ZK simulation had a gap, since fixed; see Sefranek,
``How (Not) to Simulate PLONK'' (SCN 2024).
\end{itemize}

This reflects the practical reality of deployed SNARKs: tight reductions
in idealised models, loose reductions in standard models, and deployment
selecting the former.

\section{PCD, IVC, and accumulation schemes}\label{sec:pcd}

We now turn to the most modern part of the foundations: recursive proof
composition. This is what makes Halo 2's ``no trusted setup with recursion''
combination possible, and what underlies any rollup architecture.

\subsection{Incrementally-verifiable computation}

\begin{definition}[IVC, informal]\label{def:ivc}
An \emph{incrementally-verifiable computation} (IVC; Valiant, TCC 2008) is
a scheme for proving the correctness of a computation
$z_n = f^{(n)}(z_0) := f(f(\dots f(z_0) \dots))$ such that at each step
$z_{i+1} = f(z_i)$, the prover produces a proof $\pi_{i+1}$ asserting that
$z_{i+1}$ is the correct $i+1$-st iterate. The prover at step $i+1$ has
access to $\pi_i$ and produces $\pi_{i+1}$ in time independent of $i$ (up
to logarithmic factors).
\end{definition}

The classical approach to IVC is \emph{recursive SNARK composition}: at step
$i + 1$, prove
``$z_{i+1} = f(z_i)$ AND there exists $\pi_i$ accepting against the
verifier circuit for step $i$.'' The proof $\pi_{i+1}$ thereby attests to
the full chain of steps, recursively.

For this to work, the SNARK's verifier must be efficiently expressible as
a circuit. This is non-trivial: pairing-based SNARKs require an in-circuit
pairing check (expensive), while STARK verifiers require in-circuit hash
computations (also expensive). The Halo line addresses this differently.

\begin{definition}[PCD, informal]\label{def:pcd}
\emph{Proof-carrying data} (PCD; Chiesa, Tromer; ICS 2010) generalises
IVC: instead of a linear chain of computations, the protocol supports an
arbitrary DAG of computations, each of which may consume multiple parent
proofs and produce one descendant proof.
\end{definition}

A rollup chain of $N$ transactions is naturally a PCD instance: each block
consumes the previous block's proof plus the new transaction batch, and
produces a single proof attesting to the entire history.
\subsection{Naive recursion and its failure to scale}

The naive approach to IVC proves, at step $i$, the statement ``$z_i = f(z_{i-1})$ AND
$\mathsf{Verify}(\mathsf{vk}, z_{i-1}, \pi_{i-1}) = 1$,'' compiled to a SNARK.

This succeeds in principle but encounters two practical obstacles:
\begin{enumerate}[leftmargin=*]
\item For pairing-based SNARKs, the verifier circuit contains a pairing
equation, which is costly to express algebraically inside a circuit
(it requires extension-field arithmetic).
\item For SNARKs with linear-time verifiers (such as Halo 2's $O(d)$ MSM),
the linear-time check dominates the verifier circuit, which negates
the per-step efficiency.
\end{enumerate}
The accumulation paradigm addresses (2). For (1), Halo / Halo 2 avoids
the issue by using IPA instead of pairings.

\subsection{Accumulation schemes}

The central insight, due to Bowe--Grigg--Hopwood for Halo and which BCMS20
formalised, is the following.

\begin{definition}[Accumulation scheme, informal]\label{def:acc}
For a predicate $\Phi$ (for example, ``this IPA opening is correct''), an
\emph{accumulation scheme} consists of:
\begin{itemize}[leftmargin=*]
\item An \emph{accumulator}, a compact data structure representing a batch
of claims about $\Phi$.
\item A \emph{folding} algorithm: given an old accumulator and a new claim
about $\Phi$, it produces a new accumulator that morally ``contains'' both.
The folding is fast (sub-linear in the claim size).
\item A \emph{decider}, which is expensive (linear-time) but which the scheme
invokes only once at
the very end to validate the accumulator. If the decider accepts, every
claim folded in so far is true.
\end{itemize}
\end{definition}

\begin{theorem}[BCMS20 accumulation, informal]\label{thm:bcms20}
For the IPA-PCS, there exists an accumulation scheme in which:
\begin{itemize}[leftmargin=*]
\item An accumulator is a tuple
$\mathsf{Acc} := (G^{(0)}, u_1, \dots, u_k) \in \G \times \F^k$,
representing the claim that $G^{(0)} = \Comm(g(X; u_1, \dots, u_k))$.
\item Folding two accumulators runs in $O(\log d)$ time, producing a new
accumulator that is valid iff both inputs were.
\item The decider performs one $O(d)$-size MSM to check validity.
\item Soundness reduces to DL on $\G$ in the random oracle model.
\end{itemize}
\end{theorem}

This resolves the IVC problem in the IPA setting: the in-circuit verifier
need only check the $O(\log d)$ ``succinct'' parts of an IPA proof and fold the
$O(d)$ ``deferred'' MSM check into a running accumulator. The decider pays the
deferred
work once, at the end of the entire chain.
Theorem~\ref{thm:bcms20-acc} restates this concretely once we have the IPA
construction in hand.

\subsection{The necessity of a cycle of curves}

The accumulation scheme's in-circuit folding involves arithmetic on group
elements of the underlying curve. For this to be efficient inside the
circuit, the group's base field must equal the circuit's native field.
Since the circuit's native field is the scalar field of the curve on which
the prover commits the SNARK --- which in turn is generally different from the
curve's own base field --- a difficulty arises: an $E_p$-IPA proof's group
elements have coordinates in $\F_p$, so verifying it natively would require
a circuit whose native field is $\F_p$. But a SNARK's verifier circuit
naturally operates over the curve's \emph{scalar} field $\F_q$ (where
$q = |E_p|$): the Fiat--Shamir challenges and folding scalars are $\F_q$
elements. Since a single curve cannot have $\F_p = \F_q$, no one curve
$E_p$ can host both the proof's coordinates and its own verifier circuit
natively.

The resolution is a \emph{2-cycle} of elliptic curves, where $|E_p| = q$ and
$|E_q| = p$. Then:
\begin{itemize}[leftmargin=*]
\item A circuit committed on $E_q$ verifies a proof generated over $E_p$;
that circuit's native field is the scalar field of $E_q$, namely $\F_p$
(since $|E_q| = p$) --- which coincides with the coordinate field of
$E_p$, enabling native manipulation of $E_p$'s group elements.
\item A circuit over $\F_q$ verifies the next proof, generated over $E_q$.
\item Proofs alternate between the two curves.
\end{itemize}
This is precisely the Pasta cycle of the \emph{Math Guide},
where the two curves are Pallas (over $\F_{p_{\Pallas}}$, order $p_{\Vesta}$)
and Vesta (over $\F_{p_{\Vesta}}$, order $p_{\Pallas}$).

\subsection{Connection to the Halo 2 protocol}

The Halo 2 protocol (\S\ref{sec:halo2-full}) ends each proof
with a deferred-MSM relation $G^{(0)} \stackrel{?}{=} \Comm(g)$. In a
non-recursive deployment (the current Orchard), the verifier pays this MSM
as part of the standard verification cost. In a recursive
deployment, an accumulator absorbs the same MSM check, defers it
across a chain of proofs, and the final decider pays it alone.

The Pasta cycle renders the folding step efficient. In a recursive
deployment, a Vesta-IPA proof contains in its circuit a verifier of the
previous Pallas-IPA proof. The Vesta circuit's native field is
$\F_{p_{\Pallas}}$ (the Vesta scalar field), which coincides with the
coordinate field of Pallas points. The in-circuit verifier can therefore
manipulate Pallas group elements ($\cmsym$, $\rk$, $\cvsym^{\mathsf{net}}$,
the accumulator's $G^{(0)}$) using native field arithmetic. The next
proof, a Pallas-IPA proof with native field $\F_{p_{\Vesta}}$, verifies
that Vesta proof using the symmetric argument. This ``cross-curve hop''
between Pallas and Vesta is what enables proof-carrying data without
expensive non-native field simulation.

\section{The Halo 2 proof system}\label{sec:halo2}

We now assemble the abstractions of the \emph{Crypto Guide}--\S\ref{sec:pcd} into
the concrete proof system Halo 2: a transparent PLONKish SNARK whose
polynomial commitment is the inner-product argument over the Pasta cycle.
This section develops Halo 2 to the level of detail required to state
\emph{what an Orchard Action proof actually is}, the subject of
the \emph{Orchard Guide}.

\subsection{PLONKish arithmetisation}\label{sec:plonkish}

\begin{definition}[PLONKish constraint system]\label{def:plonkish}
Fix a prime field $\F$ and an integer $k$ defining a multiplicative
subgroup $H := \{\omega^0, \dots, \omega^{n-1}\} \subset \F^*$ of size
$n := 2^k$, where $\omega$ is a primitive $n$-th root of unity. A
\emph{PLONKish circuit} consists of:
\begin{itemize}[leftmargin=*]
\item Columns partitioned into \emph{fixed} (preprocessed),
\emph{advice} (prover-witness), and \emph{instance} (public-input).
\item A maximum constraint degree $d$.
\item A finite sequence of \emph{custom gates}
$g_t(\vec x_0, \vec x_1, \dots) = 0$ --- multivariate polynomials of
degree $\le d$ in column values at the current row and at offsets
$\omega^r X$ for fixed small $r$.
\item A subset $S \subseteq H \times [\#\text{cols}]$ of
\emph{equality-constraint cells} controlled by a permutation $\sigma$.
\item A set of \emph{lookup arguments}: input tuples $(A_0, \dots, A_{m-1})$
and table tuples $(S_0, \dots, S_{m-1})$, with the assertion that for
every row $i \in [0, n)$, $(A_0(\omega^i), \dots, A_{m-1}(\omega^i))$
matches some row of the table.
\end{itemize}
\end{definition}

Vanilla PLONK (Gabizon--Williamson--Ciobotaru, 2019) employs a fixed five-selector
universal gate
\[
q_M(X) a(X) b(X) + q_L(X) a(X) + q_R(X) b(X) + q_O(X) c(X) + q_C(X) \equiv 0 \pmod{Z_H(X)},
\]
where $Z_H(X) := X^n - 1$ is the \emph{vanishing polynomial} of the domain
$H$ (it is zero exactly on $H$, so ``$\equiv 0 \bmod Z_H$'' means ``zero at
every row''). ``PLONKish'' generalises this by permitting designer-chosen
gates of higher arity and degree, each gated by a selector $q_i(X)$
vanishing on rows where the gate is inactive. The Orchard Action circuit
uses ten advice columns plus a small fixed-column lookup table for
Sinsemilla; its custom gates encode the algebraic constraints of the
Action statement (the \emph{Orchard Guide}) together with the
incomplete-addition arithmetic for in-circuit elliptic-curve operations.

\paragraph{Polynomial representation.}
The Lagrange basis $\ell_j(X)$ for $H$, where $\ell_j(\omega^j) = 1$
and $\ell_j(\omega^{j'}) = 0$ for $j' \ne j$, interpolates each column into
a polynomial of degree $< n$ over $H$. The prover commits to each
polynomial via the IPA-PCS (\S\ref{sec:ipa-detail}). All subsequent
arguments are \emph{polynomial identities} on these committed polynomials,
verifiable by querying them at random points.

\subsection{The permutation (equality) argument}\label{sec:perm}

A PLONK-style permutation argument, which Halo 2 generalises to multiple
columns, enforces equality constraints between distinct cells. Each
cell $(i, j)$ (column $i$, row $j$) receives a unique label
$\delta_i \cdot \omega^j$ with $\delta_i$ chosen so all products are
distinct as $(i, j)$ ranges over the trace. The prover writes
$\sigma$-permuted polynomials $s_i(X)$ with $s_i(\omega^j) =
\delta_{i'} \cdot \omega^{j'}$ when $\sigma(i, j) = (i', j')$.

Given Fiat--Shamir challenges $\beta, \gamma \in \F$, the prover commits
to a running-product polynomial $Z_{\mathrm{perm}}(X)$ asserting
\[
\prod_{i, j} \frac{v_i(\omega^j) + \beta\,\delta_i\,\omega^j + \gamma}
{v_i(\omega^j) + \beta\,s_i(\omega^j) + \gamma} = 1.
\]
Two values appear in equality-constrained cells iff this product over the
trace equals $1$. The verifier checks this by opening $Z_{\mathrm{perm}}$
at $x$ and $\omega x$, verifying the boundary $Z_{\mathrm{perm}}(\omega^0) = 1$
and the row-by-row update. The per-proof verifier work is $O(m + \log n)$.

\subsection{The lookup argument}\label{sec:lookup}

Halo 2's lookup argument is structurally simpler than plookup
(Gabizon--Williamson 2020) and supports arbitrary (possibly non-fixed)
tables.

\paragraph{Objective.} Establish that every entry of column $A(X)$ on the active
domain appears as some entry of column $S(X)$.

\paragraph{Construction.}
The prover supplies polynomials $A'(X), S'(X)$ where (i) $A'$ is a
permutation of $A$ with like values grouped into contiguous runs, and
(ii) $S'$ is a permutation of $S$ in which the first row of each run of
like values in $A'$ has the matching value in $S'$. The prover commits to
a running-product column $Z_{\mathrm{lookup}}(X)$, then after challenges
$\beta, \gamma$ proves:
\begin{align*}
\ell_0(X) \bigl(1 - Z_{\mathrm{lookup}}(X)\bigr) &= 0,\\
Z_{\mathrm{lookup}}(\omega X)(A'(X) + \beta)(S'(X) + \gamma) - Z_{\mathrm{lookup}}(X)(A(X) + \beta)(S(X) + \gamma) &= 0,\\
\ell_0(X)(A'(X) - S'(X)) &= 0,\\
(A'(X) - S'(X))(A'(X) - A'(\omega^{-1} X)) &= 0.
\end{align*}
The first three force $A'$ to be a permutation of $A$ and to start each
run aligned with $S'$. The fourth constrains every subsequent row to
either start a new run ($A'_i = S'_i$) or continue the previous one
($A'_i = A'_{i-1}$). Inductively, every value of $A$ equals some value of
$S$.

A random-linear combination with a challenge $\theta$ reduces multi-column
lookups to single-column:
$A_{\mathrm{c}}(X) := \sum_j \theta^{m-1-j} A_j(X)$. Variable tables
(where $S$ is itself an advice column) and range checks
(where $S$ enumerates the allowed range) are special cases.

\subsection{The inner-product polynomial commitment}\label{sec:ipa-detail}

The Halo 2 polynomial commitment, instantiating the abstract PCS of
the \emph{Crypto Guide}, is the inner-product argument of
Bootle--Cerulli--Chaidos--Groth--Petit (EUROCRYPT 2016) as refined in
Bulletproofs (B\"unz et al., 2017). We give the construction in detail.
Write $d := 2^k$ for the degree bound (equivalently, the number of
generators); in the assembled protocol of \S\ref{sec:halo2-full} the
committed polynomials have degree $< n$, so $d = n$ there.

\begin{construction}[IPA-based polynomial commitment]\label{const:ipa}
Let $\G$ be a cyclic group of prime order $p$ in which DL is hard. Choose
$d = 2^k$ and sample, via hash-to-curve, $\vec{G} = (G_1, \dots, G_d)$
and $H \in \G$ with no known non-trivial linear relation among them.

For $p(X) = \sum_{i=0}^{d-1} a_i X^i \in \F_p[X]$ and blinder
$r \in \F_p$,
\[
\Comm(p; r) := \ip{\vec{a}}{\vec{G}} + [r] H \in \G.
\]
\end{construction}

To open $p$ at $x \in \F_p$ with claimed value $v$, the relation is
\[
\Rel_{\mathrm{IPA}} := \Bigl\{\bigl((P, x, v); (\vec{a}, r)\bigr) :
P = \ip{\vec{a}}{\vec{G}} + [r] H,\;
v = \ip{\vec{a}}{\vec{b}(x)},\; \deg p \le d-1\Bigr\},
\]
with $\vec{b}(x) := (1, x, x^2, \dots, x^{d-1})$.

\subsubsection{The recursive halving protocol}

The argument runs in $k = \log_2 d$ rounds. Initialise
$\vec{a}^{(k)} := \vec{a}$, $\vec{G}^{(k)} := \vec{G}$,
$\vec{b}^{(k)} := \vec{b}(x)$, and let $U \in \G$ be a challenge group
element (in Fiat--Shamir, derived from the transcript). The prover
absorbs $[v]U$ so that the relation becomes
\[
P + [v] U = \ip{\vec{a}}{\vec{G}} + [r] H + [\ip{\vec{a}}{\vec{b}}] U.
\]

In round $j$ (descending from $k$ to $1$): split each vector into halves
and let
\begin{align*}
L_j &:= \ip{\vec{a}_{\mathrm{lo}}^{(j)}}{\vec{G}_{\mathrm{hi}}^{(j)}} + [\ell_j] H + [\ip{\vec{a}_{\mathrm{lo}}^{(j)}}{\vec{b}_{\mathrm{hi}}^{(j)}}] U,\\
R_j &:= \ip{\vec{a}_{\mathrm{hi}}^{(j)}}{\vec{G}_{\mathrm{lo}}^{(j)}} + [r_j] H + [\ip{\vec{a}_{\mathrm{hi}}^{(j)}}{\vec{b}_{\mathrm{lo}}^{(j)}}] U.
\end{align*}
The verifier returns a uniform $u_j \in \F_p^*$. Both parties fold:
\begin{align*}
\vec{a}^{(j-1)} &:= u_j\,\vec{a}_{\mathrm{lo}}^{(j)} + u_j^{-1}\,\vec{a}_{\mathrm{hi}}^{(j)},\\
\vec{G}^{(j-1)} &:= u_j^{-1}\,\vec{G}_{\mathrm{lo}}^{(j)} + u_j\,\vec{G}_{\mathrm{hi}}^{(j)},\\
\vec{b}^{(j-1)} &:= u_j^{-1}\,\vec{b}_{\mathrm{lo}}^{(j)} + u_j\,\vec{b}_{\mathrm{hi}}^{(j)}.
\end{align*}

After $k$ rounds each vector has length 1. The prover sends the final
scalar $a := \vec{a}^{(0)}$ and an aggregated blinder $r^*$. The verifier
accepts iff
\begin{equation}\label{eq:ipa-final}
P + [v] U + \sum_{j=1}^{k} \bigl([u_j^2] L_j + [u_j^{-2}] R_j\bigr)
\;\stackrel{?}{=}\; [a]\,G^{(0)} + [r^*] H + [a \cdot b^{(0)}] U,
\end{equation}
where $G^{(0)} = \ip{\vec{s}}{\vec{G}}$ and the challenges determine
$\vec{s}$ as below.

\subsubsection{The structured scalars and the polynomial $g$}\label{sec:gpoly}

Write the binary expansion $i = \sum_{\ell=0}^{k-1} b_\ell 2^\ell$ for
$i \in [0, d)$. By induction on the folding,
\[
G^{(0)} = \ip{\vec{s}}{\vec{G}}, \qquad b^{(0)} = \ip{\vec{s}}{\vec{b}(x)} = g(x; u_1, \dots, u_k),
\]
where
\[
s_i := \prod_{\ell=0}^{k-1} u_{\ell+1}^{(-1)^{1-b_\ell}}, \qquad
g(X; u_1, \dots, u_k) := \prod_{\ell=1}^{k} \bigl(u_\ell^{-1} + u_\ell X^{2^{\ell-1}}\bigr).
\]

The essential asymmetry is the following: $g(X)$ has degree $d - 1$ and the verifier can evaluate
it at $x$ in $O(\log d)$ field operations; but computing
$G^{(0)} = \Comm(g)$ requires an $O(d)$-size multi-scalar
multiplication (MSM). This asymmetry is the seed of the accumulation
scheme (\S\ref{sec:halo-acc}).

\subsubsection{Worked example: an IPA opening at \texorpdfstring{$d = 4$}{d=4}}\label{sec:ipa-worked}

Take $d = 4$, so $k = 2$ rounds. Let $p(X) = a_0 + a_1 X + a_2 X^2 + a_3 X^3$
with $\vec{a} = (a_0, a_1, a_2, a_3)$. The verifier holds
$P = \ip{\vec{a}}{\vec{G}} + [r] H$ (we suppress $r$ to declutter; the
$H$-part propagates identically) and requires $v = p(x)$ at some $x \in \F_p$.

Set $\vec{b}^{(2)} := (1, x, x^2, x^3)$ and $\vec{G}^{(2)} := \vec{G}$.
After the prover absorbs $[v] U$, the relation is
$P + [v] U = \ip{\vec{a}}{\vec{G}} + \ip{\vec{a}}{\vec{b}^{(2)}}\,U$.

\paragraph{Round 2.} Split into halves of size 2. The prover sends
\begin{align*}
L_2 &= a_0 G_3 + a_1 G_4 + (a_0 x^2 + a_1 x^3) U,\\
R_2 &= a_2 G_1 + a_3 G_2 + (a_2 + a_3 x) U.
\end{align*}
The verifier picks $u_2 \in \F_p^*$ and both parties fold to length-2
vectors as above.

\paragraph{Round 1.} Split the length-2 vectors. The prover sends $L_1,
R_1$ analogously; the verifier picks $u_1$. After folding, all three vectors
have length 1.

\paragraph{Final check.} With $i = b_0 + 2 b_1$ ranging over
$\{0, 1, 2, 3\}$:
\begin{align*}
s_0 &= u_1^{-1} u_2^{-1}, & s_1 &= u_1 u_2^{-1}, & s_2 &= u_1^{-1} u_2, & s_3 &= u_1 u_2.
\end{align*}
The verifier evaluates
$g(x; u_1, u_2) = (u_1^{-1} + u_1 x)(u_2^{-1} + u_2 x^2)$ in two
multiplications, and computes
$G^{(0)} = \sum_{i=0}^{3} [s_i] G_{i+1}$ as a length-4 MSM. The check
(\ref{eq:ipa-final}) becomes
\[
P + [v] U + [u_1^2] L_1 + [u_1^{-2}] R_1 + [u_2^2] L_2 + [u_2^{-2}] R_2
\;\stackrel{?}{=}\; [a] G^{(0)} + [a \cdot g(x; u_1, u_2)]\,U.
\]

The folding rule for $\vec{G}$ ensures that $G^{(0)}$ is a specific
linear combination of the original $\vec{G}$ with coefficients $\vec{s}$;
that the $s_i$ match $\prod_\ell (u_\ell^{-1} \text{ or } u_\ell)$
according to the binary expansion of $i$ is exactly the closed-form
identity. The verifier can evaluate $g$ in $\log d$ field operations but
recomputing $G^{(0)}$ takes $O(d)$ MSM --- the asymmetry the accumulator
defers.

\subsubsection{Knowledge soundness via rewinding}\label{sec:ipa-extract}

\begin{theorem}[Knowledge soundness of IPA]\label{thm:ipa-ks}
Under the DL assumption on $\G$, and modelling Fiat--Shamir as a random
oracle, Construction~\ref{const:ipa} is knowledge-sound for
$\Rel_{\mathrm{IPA}}$ with knowledge error $O(d / |\F_p|)$.
\end{theorem}

\begin{proof}[Sketch of the tree extractor]
The extractor rewinds the prover $\mathcal{P}^*$ to obtain, at each
round, two transcripts that agree on $(L_j, R_j)$ but differ in $u_j$.
The two corresponding verifier equations form a $2 \times 2$ linear
system in the round-$j$ halved witness whose Vandermonde determinant is
$(u_j^2 - u_j'^2)/(u_j u_j')$, non-zero whenever $u_j \ne \pm u_j'$. This
recovers the witness halves $\vec{a}_{\mathrm{lo}}^{(j)},
\vec{a}_{\mathrm{hi}}^{(j)}$.

Extracting the full witness $\vec{a} \in \F_p^d$ requires two successful
transcripts at the deepest level, then for \emph{each pair}, two
transcripts at the next-shallower level, and so on. The tree has
$2^k = d$ leaves, yielding polynomial expected extraction time when
$d = \poly(\lambda)$. The knowledge error of $O(d / |\F_p|)$ accounts
for the probability that any $2d$ challenges in the tree happen to
collide, which a union bound renders negligible.

This is the recursive generalisation of 2-special soundness for Schnorr
(the \emph{Crypto Guide}): a single transcript reveals nothing, but two
transcripts on the same first message permit inversion of a $2 \times 2$
system; in IPA the inversion is structured but the principle is identical.
\end{proof}
\subsection{Accumulation and the Halo trick}\label{sec:halo-acc}

The verifier's check (\ref{eq:ipa-final}) runs in $O(\log d)$ scalar
operations and $\log d$ point additions, \emph{except} for the single
evaluation $G^{(0)} = \ip{\vec{s}}{\vec{G}}$ --- an MSM of length $d$.
The accumulation idea (Bowe--Grigg--Hopwood 2019; formalised by BCMS20)
is to \emph{defer} this MSM and amortise it across many proofs.

\begin{definition}[Atomic accumulator for IPA]\label{def:ipa-acc}
An IPA \emph{accumulator instance} is a tuple
$\Acc := (G^{(0)}, u_1, \dots, u_k) \in \G \times \F_p^k$. It is
\emph{valid} iff $G^{(0)} = \Comm(g(X; u_1, \dots, u_k))$.
\end{definition}

\begin{theorem}[Accumulation of IPA, BCMS20]\label{thm:bcms20-acc}
There is an accumulation scheme
$(\mathsf{AccVerifier}, \mathsf{AccDecider})$ for IPA with:
\begin{itemize}[leftmargin=*]
\item $\mathsf{AccVerifier}$, on input a new IPA instance and an old
accumulator, runs in $O(\log d)$ and outputs a new accumulator $\Acc'$.
\item $\mathsf{AccDecider}$, run \emph{once} at the end of a chain, performs
an MSM of length $d$ to check $G^{(0)} = \Comm(g)$.
\item Soundness reduces to DL on $\G$ in the random-oracle model.
\end{itemize}
\end{theorem}

\paragraph{Operational interpretation.}
A Halo 2 verifier circuit, expressed over the scalar field of one Pasta
curve and verifying a proof produced over the \emph{other} curve, can:
\begin{enumerate}[leftmargin=*]
\item Run the $O(\log d)$ ``succinct'' checks of the IPA verifier natively.
\item Witness the next-step claimed $G^{(0)\prime}$ and the next challenges
$u_1', \dots, u_k'$ as public inputs to the next proof, deferring the MSM
check.
\item Fold the deferred MSM checks: by linear combination with a random
challenge $\rho$, reduce two MSM checks to one,
$\Acc'' := \Acc + \rho \cdot \Acc'$.
\end{enumerate}
A single final MSM at the end of the chain validates the entire history.
The Pasta cycle ensures the in-circuit scalar arithmetic on each step is
native, as we develop in \S\ref{sec:pcd}.

The mathematical core is the recognition that
$G^{(0)} = \Comm(g(X; u_1, \dots, u_k))$ --- the linear-time check is
\emph{itself a polynomial commitment opening of a structured polynomial}
and hence amenable to recursion within an IPA.

\subsection{The complete Halo 2 protocol}\label{sec:halo2-full}

We now assemble the pieces into the seven-phase protocol implemented in
the halo2 codebase. Throughout, $\F$ denotes the constraint field
(Orchard: $\F_{p_{\Pallas}}$), $\G$ the IPA group (Orchard:
$\Vesta(\F_{p_{\Vesta}})$), and $H$ the multiplicative subgroup of $\F$
of order $n$.

\paragraph{Setup (preprocessed, one-time).}
The circuit designer specifies columns, custom gates, lookup tables, and
the equality permutation $\sigma$. Both parties derive commitments to the
fixed-column polynomials $f_i(X)$, selector polynomials $q_i(X)$,
permutation polynomials $s_i(X)$, and lookup-table polynomials $S_j(X)$,
together with the Lagrange basis on $H$, the vanishing polynomial
$Z_H(X) := X^n - 1$, and the IPA generators $\vec{G} \in \G^n$ and
$H \in \G$.

\paragraph{Phase 1 --- advice commitments.}
The prover constructs the advice column values $\vec{w_i} \in \F^n$
satisfying all gates and equality constraints; interpolates each
$\vec{w_i}$ to a polynomial $w_i(X)$ of degree $< n$ on $H$; adds a
random low-degree blinding polynomial for ZK; commits via IPA-PCS; sends
the commitments.

\paragraph{Phase 2 --- lookup challenges $\theta, \beta, \gamma$.}
For each lookup, the verifier (via Fiat--Shamir) returns challenges:
$\theta$ to compress multi-column lookups, $\beta, \gamma$ to instantiate
the running-product. The prover commits to $A', S', Z_{\mathrm{lookup}}$.

\paragraph{Phase 3 --- permutation argument.}
The prover constructs $Z_{\mathrm{perm}}(X)$ as in \S\ref{sec:perm} and
commits to it.

\paragraph{Phase 4 --- quotient challenge $y$.}
The verifier returns a fresh $y \in \F$. The prover assembles the
\emph{gate equation}: a single polynomial $C(X)$ that is a $y$-weighted
linear combination of (i) each custom gate, (ii) permutation constraints,
(iii) lookup constraints. By construction $C(X)$ vanishes on $H$ when
all gates pass. The prover computes the quotient
$h(X) := C(X) / Z_H(X) = C(X) / (X^n - 1)$, splits $h$ into chunks of
degree $< n$, and commits to each chunk.

\paragraph{Phase 5 --- evaluation challenge $x$.}
The verifier returns $x \in \F$. The prover evaluates every committed
polynomial at $x$, and at $\omega x$ for adjacent-row references, and
sends these claimed evaluations as field elements. The verifier checks
the gate equation $C(x) = h(x) Z_H(x)$ as a numerical identity in the
sent evaluations. This check is meaningful only once the claimed
evaluations bind to the commitments, which is the role of phases 6--7;
a prover who sent inconsistent evaluations would pass this field identity
but fail the opening argument.

\paragraph{Phase 6 --- batched polynomial opening.}
To verify the evaluations are consistent with the commitments, the prover
runs Halo 2's multipoint-opening argument: it groups the polynomials queried at the
same point set and combines them via a challenge $x_1$, then builds a
single polynomial $f^*(X)$ representing all openings such that
opening $f^*$ at one final point $x_3$ implies all original openings.

\paragraph{Phase 7 --- IPA opening.}
Prover and verifier execute the recursive halving IPA
(\S\ref{sec:ipa-detail}) on $f^*$, yielding $\log_2 n$ pairs $(L_j, R_j)$,
a final scalar $a$, and a final blinder $r^*$. The verifier accepts iff
(\ref{eq:ipa-final}) holds.

\paragraph{Accumulation (recursion).}
The next proof's verifier circuit \emph{defers} the MSM check
$G^{(0)} = \Comm(g(X; u_1, \dots, u_k))$ at the end of
phase 7 into an accumulator and folds it with prior
accumulators (\S\ref{sec:halo-acc}).
A single decider at the very end validates the entire chain.

\paragraph{Costs.}
Per proof, ignoring the deferred MSM: $O(\log n)$ pairs checked in
phase 7; one gate-equation field check; a constant number of polynomial
evaluations consumed. Without accumulation, the verifier additionally
performs one $O(n)$-size MSM; with accumulation the chain pays for this MSM once at
the end. FFTs and one $O(n)$ MSM dominate the prover's work
(typical Orchard Action: $n = 2^{11}$, a few seconds on commodity
hardware).


\subsection{Instance columns and binding the public statement}\label{sec:instance}

Definition~\ref{def:plonkish} lists \emph{instance} columns alongside fixed
and advice columns, but the arguments of \S\ref{sec:perm}--\S\ref{sec:halo2-full}
never single them out. We single them out now, because what an Orchard Action proof
\emph{asserts} is precisely a statement about its instance values --- the
anchor, nullifier $\nfsym$, net value commitment $\cvsym^{\mathsf{net}}$,
randomised key $\rk$, and the $\mathsf{enableSpends}/\mathsf{enableOutputs}$
flags (the \emph{Orchard Guide}).

\paragraph{The instance polynomial.}
An instance column is a public-input column: its $n$ entries
$\vec{\pi} = (\pi_0, \dots, \pi_{n-1}) \in \F^n$ are not part of the
witness but both parties agree on them before verification. Like every
other column it is interpolated over $H$ in the Lagrange basis,
\[
\pi(X) := \sum_{j=0}^{n-1} \pi_j\, \ell_j(X) \in \F[X], \qquad \deg \pi < n,
\]
where $\ell_j(\omega^j) = 1$ and $\ell_j(\omega^{j'}) = 0$ otherwise
(\S\ref{sec:plonkish}). The prover writes the public values into a handful of
designated rows; all remaining $\pi_j$ are $0$. A custom gate of the form
$q(X)\bigl(a(X) - \pi(X)\bigr) = 0$, with $q$ a selector active exactly on
those rows, copies an advice cell to a public value and thereby asserts
``advice cell $= \pi_j$.'' In Orchard, the gate constrains the cells carrying the in-circuit
$\nfsym$, $\cvsym^{\mathsf{net}}$, $\rk$ and the flags
against the instance column this way, so a satisfying assignment exists
\emph{only} for the public values the verifier supplied.

\paragraph{Where it enters the protocol.}
Unlike advice, the prover never commits to $\pi(X)$ and never includes it
among the IPA openings of phases~6--7: the verifier already holds
$\vec{\pi}$, so committing would be redundant and would disclose nothing. The
instance polynomial enters only the \emph{gate equation}. When the
verifier checks $C(x) = h(x)\,Z_H(x)$ at the evaluation challenge $x$
(Phase~5, \S\ref{sec:halo2-full}), the assembled $C(x)$ contains the term
$q(x)\bigl(a(x) - \pi(x)\bigr)$ from the gate above. The advice evaluation
$a(x)$ arrives in the transcript and the opening argument binds it to its commitment; the
selector evaluation $q(x)$ is a fixed-column value
both parties recompute. The verifier supplies the remaining evaluation
$\pi(x)$ itself, computing $\pi(x) = \sum_j \pi_j\, \ell_j(x)$ directly
from the known public values --- an $O(\log n)$ evaluation using the
barycentric form of the Lagrange basis on $H$, never read from the proof.

\paragraph{Consequence.}
A prover who substitutes different public values $\vec{\pi}\,'$ must still
satisfy $C(x) = h(x)\,Z_H(x)$ for the $\pi(x)$ the \emph{verifier}
computes from the genuine $\vec{\pi}$. Knowledge soundness
(\S\ref{sec:ks}) then extracts a witness for the relation instantiated at
exactly those public values; no witness exists for the wrong ones except
with negligible probability. This binds the proof to its public
statement: an Action proof that a node verifies against an anchor it accepts,
a nullifier it will mark spent, and a value commitment it will sum into
the net balance check (the \emph{Orchard Guide}) can only verify if the
prover knew a witness consistent with \emph{those} values.

\subsection{Zero knowledge: hiding the witness in Halo 2}\label{sec:halo-zk}

Knowledge soundness (\S\ref{sec:ks}) establishes that a convincing prover \emph{knows} a
witness; it remains to explain why the proof \emph{reveals} no more than that. Halo~2 is
zero-knowledge through two devices --- hiding commitments and blinded polynomials
--- assembled into a simulator.

\paragraph{Hiding commitments.} The polynomial commitment of
\S\ref{sec:ipa-detail} is \emph{perfectly hiding}. The prover commits a polynomial \(p\)
as
\[
  \mathrm{Commit}(p; r) \;=\; \ip{\vec{p}}{\vec{G}} + [r]\,H, \qquad
  r \stackrel{\$}{\leftarrow} \F_p,
\]
with \(\vec p\) the coefficient vector of \(p\) and \(H\) an independent generator
of no known discrete-log relation to \(\vec G\). The \([r]H\) term makes the
commitment a \emph{uniform} group element for every \(p\), so a commitment by
itself leaks nothing about what the prover committed (the perfect hiding of a Pedersen
commitment, \emph{Crypto Guide}). Every commitment the verifier sees --- the
advice columns, the permutation and lookup products, the quotient chunks --- has
this form.

\paragraph{Blinding the witness polynomials.} The prover eventually
\emph{opens} a commitment, and an opening reveals evaluations. The prover interpolates the advice columns
over the domain \(H\) to polynomials of degree \(< n\); to prevent those
evaluations from leaking the witness, the prover does not use all \(n\) rows for
circuit data. It reserves a small fixed number of rows at the bottom of the table
--- the \emph{blinding rows}, never covered by an active gate --- and fills them
with uniform random field elements, at least one more than the number of times the
verifier will open any single column. Its \(n\) evaluations pin down a degree-\((<n)\) polynomial,
so with enough independent random values planted on the
blinding rows the handful the verifier actually queries are jointly uniform and
independent of the circuit data. The prover blinds the permutation and lookup running products and
the quotient pieces the same way.

\paragraph{Why the openings leak nothing.} In each phase the verifier sees only
(i) hiding commitments --- uniform group elements --- and (ii) a bounded number of
evaluations of the committed polynomials at the Fiat--Shamir challenge points. The
blinding rows make each such evaluation uniformly random subject only to the
\emph{public} relations the verifier itself checks; no information about the
private witness passes through. The final inner-product opening
(\S\ref{sec:ipa-detail}) likewise runs with a blinder, so it too hides its
input.

\paragraph{The simulator.} Formally, zero knowledge is the existence of a
\emph{simulator} that, given only the public input and \emph{not} the witness,
outputs a transcript indistinguishable from a real one. Halo~2's simulator runs
the protocol ``backwards,'' exactly as Schnorr's does (\emph{Crypto Guide}): it
chooses the evaluations and the final opening so as to satisfy the verifier's
equations, then chooses hiding commitments consistent with them --- possible
precisely because the commitments are perfectly hiding and it may choose the blinders
freely --- and, non-interactively, \emph{programs} the Fiat--Shamir oracle so
the challenges come out as required. Since the commitments are perfectly hiding
and a blinder protects every opening, the simulated transcript is distributed as a real
one. Interactively this is honest-verifier zero knowledge; compiled through
Fiat--Shamir in the random-oracle model it is a NIZK that reveals nothing beyond
the truth of the statement.

\paragraph{The three properties together.} For the relation fixed by a PLONKish
circuit, Halo~2 thus delivers all three guarantees a SNARK should:
\emph{completeness} (an honest prover holding a satisfying witness always
convinces the verifier), \emph{knowledge soundness} (\S\ref{sec:ks}: from any
convincing prover an extractor pulls a satisfying witness, so an accepted proof
certifies one is \emph{known}), and \emph{zero knowledge} (a witness-free
simulator reproduces the proof, so it leaks nothing else). It is exactly this
combination --- ``the prover knows a satisfying witness, and the proof discloses
nothing more'' --- on which the Orchard Action proof rests: knowledge soundness
yields balance and nullifier integrity, zero knowledge yields privacy
(\emph{Orchard Guide}).
\end{document}
